\section{Conclusion} \label{sec:Conclusion}

This work compared the architecture of \acp{SNN} to traditional \acp{ANN} and investigated execution methodologies on \acp{GPU}.
The first part focused on the architectural differences between the two types of neural networks.
\acp{SNN} use spike-based, time-dependent signals and are therefore more biologically plausible.
Their neuron models, signal encoding, activation mechanisms and learning rules differ from those of \acp{ANN}.
The components of \acp{SNN} vary in terms of biological plausibility and computational complexity, which necessitates trade-offs when choosing appropriate network architectures.
For example, a Hodgkin-Huxley model is more realistic than a \ac{LIF} model but significantly increases complexity.
\acp{ANN}, on the other hand, offer a more straightforward architecture with less variation.

The differences between the network types lead to distinct requirements, especially with regard to hardware compatibility, which was examined in the second part.
While \acp{GPU} are well-suited for the parallel execution of synchronous computations as used in \acp{ANN}, the asynchronous and event-driven patterns in \acp{SNN} require alternative strategies.
This work presented several approaches to address this mismatch.
These include low-level \ac{GPU} optimizations such as Temporal Fusion, as well as framework-based implementations like Brian with its \ac{GPU}-optimized extension Brian2CUDA, or Norse, which leverages PyTorch capabilities.
Benchmarks from the literature demonstrated significant speedups, particularly for large-scale networks, compared to execution on \acp{CPU}.

Although neuromorphic hardware remains more suitable for \acp{SNN} in terms of energy efficiency, the findings indicate that \acp{GPU} can serve as a practical platform for their simulation and training.
This supports further development and research in the field of \acp{SNN}.


